\documentclass[aps,pra, onecolumn,superscriptaddress,nopacs,floatfix, showkeys, notitlepage]{revtex4-1}
%   \documentclass[10pt]{article}



\usepackage{graphicx, color}
\usepackage{amssymb, amsmath}
\usepackage{times}
\usepackage[normalem]{ulem}
\usepackage{bbm}
\usepackage{todonotes}
\usepackage{mathtools}
\usepackage{enumitem} 
\usepackage{physics}
\usepackage{braket}
\usepackage{soul,xcolor}
\usepackage{subcaption}
\usepackage{bbold}
\usepackage{wrapfig}
\usepackage{multirow}

\usepackage{booktabs}  % for \toprule, \midrule, and \bottomrule macros 
\usepackage{tabularx}  % for 'tabularx' environment and 'X' column type
%\usepackage[hang,small,bf]{caption}
%\setlength{\captionmargin}{25pt}
%\usepackage{dblfloatfix}
\usepackage{tikz}
%%%<
\usepackage{verbatim}
%\usepackage[active,tightpage,floats]{preview}
%\PreviewEnvironment{tikzpicture}
%\setlength\PreviewBorder{5pt}%
%%%>
\usetikzlibrary{arrows}
\usetikzlibrary{positioning}
% Nice captions.
%\usepackage[hang,small,bf]{caption}
%\setlength{\captionmargin}{25pt}

% New commands to keep things tidy.
%\newcommand{\ket}[1]{$\left|#1\right\rangle$}
\newcommand{\Om}[1]{\small $\omega_{#1}$}
\newcommand{\De}[1]{$\Delta_{#1}$}
\newcommand{\Ga}[1]{$\Gamma_{#1}$}

%\newcommand{\Om}[1]{\small $\omega_{#1}$}
%\newcommand{\De}[1]{$\Delta_{#1}$}
%\newcommand{\Ga}[1]{$\Gamma_{#1}$}

\makeatletter
\newcommand*\bigcdot{\mathpalette\bigcdot@{.5}}
\newcommand*\bigcdot@[2]{\mathbin{\vcenter{\hbox{\scalebox{#2}{$\m@th#1\bullet$}}}}}
\newcommand*\cav{(\bigcdot)}
\newcommand*\ncav{(\bigcdot \bigcdot)}
\makeatother

%\newcommand{\ket}[1]{$\left|#1\right\rangle$}

\usepackage{amsmath}
\makeatletter
\def\env@cases{%
  \let\@ifnextchar\new@ifnextchar
  \left\lbrace
  \def\arraystretch{1.2}%
  \array{l@{\quad}l@{}}% Formerly @{}l@{\quad}l@{}
}



\begin{document}
\title{Note on Phase shift measurment of Chirped FBG}
\author{NanoQT}
\date{\today}

\begin{abstract}
\end{abstract}


\maketitle
%\tableofcontents



% \subsection{Fiber Bragg grating}
% The transmission of FBG mirror 
% \begin{figure}[!h]  
%         \centering
%         \includegraphics[width=1\linewidth]{fig/phase.png}
% 	\caption{Transmission/Reflection spectra and phase shift of a single FBG, estimated by a 1D model including the spatial inhomogenity of refractive index modulation $\delta n$. The inhomogenity is hand-picked by using a gaussian envelope with $\sigma=5$ mm and peak $\delta n = 3.6\times 10^{-4}$, leading to the asymmetry of the transmission/reflection spectra with respect to detuning (typically observed in the current system). The most right plot shows $\frac{\partial \phi_\mathrm{FBG}}{\partial}\sim-0.037$ rad/GHz at the mid gap
%  %and $\frac{\partial \phi_\mathrm{FBG}}{\partial}\sim-0.1$ rad/GHz near the bandedge.
%  }
%     \label{fig:1dsim}
% \end{figure}


%

\section{Setup}

A CW laser at optical frequency $\omega_c$ is phase-modulated by an EOM at
frequency $\Delta$ with modulation index $m$. In the slowly-varying
envelope (carrier frame),
\begin{equation}
    E_{\mathrm{in}}(t) = E_0 \, e^{\,i m \sin(2\pi \Delta t)}
    = E_0 \sum_{k=-\infty}^{\infty} J_k(m)\, e^{\,i 2\pi k \Delta t}.
\end{equation}
For small $m$ we keep only $k \in \{-1,0,+1\}$ and use
$J_{-1}(m) = -J_1(m)$:
\begin{equation}
    E_{\mathrm{in}}(t) \approx E_0\Big[
    J_0(m) + J_1(m)\,e^{\,i 2\pi \Delta t}
    - J_1(m)\,e^{-i 2\pi \Delta t}
    \Big].
\end{equation}
The carrier is at detuning $\delta$ from the FBG resonance; the two
sidebands are at $\delta \pm \Delta$.

\section{Reflected field}

The FBG is linear and time-invariant, so each spectral component acquires
its own complex reflection coefficient:
\begin{equation}
    E_r(t) = E_0\Big[
    J_0\, r(\delta)
    + J_1\, r(\delta+\Delta)\, e^{\,i 2\pi \Delta t}
    - J_1\, r(\delta-\Delta)\, e^{-i 2\pi \Delta t}
    \Big],
\end{equation}
writing $J_k \equiv J_k(m)$ and leaving the common $e^{\,i\omega_c t}$
factor implicit.

\section{Photodetector output}

A square-law detector responds to $P(t) = |E_r(t)|^2$. Let
\begin{equation}
    A = J_0\, r(\delta), \qquad
    B = J_1\, r(\delta+\Delta), \qquad
    C = -J_1\, r(\delta-\Delta),
\end{equation}
so that $E_r / E_0 = A + B\, e^{i 2\pi\Delta t} + C\, e^{-i 2\pi\Delta t}$.
Multiplying out,
\begin{equation}
    \frac{P(t)}{|E_0|^2}
    = \underbrace{|A|^2 + |B|^2 + |C|^2}_{\text{DC}}
    + \underbrace{\big(A C^* + B A^*\big)\, e^{\,i 2\pi \Delta t}
                 + \mathrm{c.c.}}_{\text{beat at } \Delta}
    + \underbrace{B C^*\, e^{\,i 4\pi \Delta t} + \mathrm{c.c.}}_{\text{beat at } 2\Delta}.
\end{equation}
We keep only the component at $\Delta$.

\section{The $\Delta$ beat}

Substituting $A$, $B$, $C$,
\begin{equation}
    A C^* + B A^*
    = -J_0 J_1\, r(\delta)\, r^*(\delta-\Delta)
      + J_0 J_1\, r(\delta+\Delta)\, r^*(\delta).
\end{equation}
Define the (normalized) complex observable
\begin{equation}
    \boxed{\;
    \mathcal{C}(\delta;\Delta)
    \;\equiv\; r^*(\delta)\,r(\delta+\Delta) \;-\; r(\delta)\,r^*(\delta-\Delta)
    \;}
\end{equation}
so that
\begin{equation}
    P(t) = P_{\mathrm{DC}} + 2\,|E_0|^2 J_0 J_1 \,
    \Re\!\big[\mathcal{C}(\delta;\Delta)\, e^{\,i 2\pi \Delta t}\big]
    + (\text{terms at } 2\Delta).
\end{equation}

\section{Demodulation}

Mixing $P(t)$ with an electrical LO at $\Delta$ and low-passing selects
one quadrature of $\mathcal{C}$:
\begin{equation}
    \big\langle P(t)\, e^{-i 2\pi \Delta t}\big\rangle_t
    = |E_0|^2 J_0(m) J_1(m)\,\mathcal{C}(\delta;\Delta).
\end{equation}
Equivalently, a Fourier transform of $P(t)$ evaluated at $+\Delta$ yields
the same $\mathcal{C}$ up to window factors.

\section{Experimental observable}

The directly measured power at frequency $\Delta$ is
\begin{equation}
    |\mathcal{C}(\delta;\Delta)|^2
    = \big|\,r^*(\delta)\,r(\delta+\Delta) - r(\delta)\,r^*(\delta-\Delta)\,\big|^2.
\end{equation}
The phase $\arg\mathcal{C}$ is obtained by I/Q demodulation.

\section{Small-$\Delta$ expansion}

Write $r(\delta) = R(\delta)\, e^{\,i\phi(\delta)}$ and Taylor expand
$r(\delta\pm\Delta) = r \pm \Delta\, r' + \tfrac{1}{2}\Delta^2\, r'' + O(\Delta^3)$.
Taking the conjugate on the $-\Delta$ side (conjugation and
$d/d\delta$ commute for real $\delta$),
\begin{align}
    r^*(\delta)\,r(\delta+\Delta)
    &= |r|^2 + \Delta\, r^* r' + \tfrac{1}{2}\Delta^2\, r^* r'' + O(\Delta^3), \\
    r(\delta)\,r^*(\delta-\Delta)
    &= |r|^2 - \Delta\, r\, r'^{\,*} + \tfrac{1}{2}\Delta^2\, r\, r''^{\,*} + O(\Delta^3).
\end{align}
Subtracting,
\begin{equation}
    \mathcal{C}(\delta;\Delta)
    = \Delta\,\big(r^* r' + r\, r'^{\,*}\big)
    + \tfrac{1}{2}\Delta^2\,\big(r^* r'' - r\, r''^{\,*}\big)
    + O(\Delta^3).
\end{equation}

Using $r^* r' = \tfrac{1}{2}\,d|r|^2/d\delta + i\,|r|^2 \phi'$, the first
bracket is purely real and the second purely imaginary:
\begin{equation}
    r^* r' + r\, r'^{\,*} = 2\,\Re(r^* r') = \frac{d|r|^2}{d\delta}, \qquad
    r^* r'' - r\, r''^{\,*} = 2i\,\Im(r^* r'').
\end{equation}
A short calculation with $r = R e^{i\phi}$ gives
$\Im(r^* r'') = 2 R R' \phi' + R^2 \phi'' = d(R^2 \phi')/d\delta$. Hence
\begin{equation}
\boxed{\;
    \mathcal{C}(\delta;\Delta)
    \;\approx\; \Delta\,\frac{d|r|^2}{d\delta}
    \;+\; i\,\Delta^2\,\frac{d\!\left(|r|^2 \phi'\right)}{d\delta}
    \;+\; O(\Delta^3).
\;}
\end{equation}

The parity $\mathcal{C}(-\Delta) = -\mathcal{C}^*(\Delta)$ forces odd
coefficients to be real and even coefficients imaginary; the expansion
above satisfies this.

\paragraph{Consequences.}
\begin{itemize}\setlength\itemsep{2pt}
  \item $\Re\mathcal{C}$ is \emph{leading order} in $\Delta$ and tracks
        $d|r|^2/d\nu$ — the amplitude derivative. Classical PDH
        error-signal shape.
  \item $\Im\mathcal{C}$ is \emph{one order smaller} and mixes $\phi'$ with
        the amplitude variation $R'$. It is not a clean proxy for
        $d\phi/d\nu$ unless $R$ is approximately constant over the scan.
\end{itemize}

\section{Phase-slope reconstruction}

The cleanest finite-$\Delta$ estimator of $d\phi/d\nu$ uses the raw phases
at the two sideband frequencies directly (both accessible once $r$ is
sampled, and experimentally via I/Q demodulation giving complex
amplitudes rather than power):
\begin{equation}
\boxed{\;
    \widehat{\phi'}(\delta)
    \;=\; \frac{\arg r(\delta+\Delta) - \arg r(\delta-\Delta)}{2\,\Delta}
    \;=\; \phi'(\delta) + \tfrac{1}{6}\Delta^2\,\phi'''(\delta) + O(\Delta^4).
\;}
\end{equation}
Any branch jumps in $\arg r$ must be unwrapped before differencing.

Using $\mathcal{C}$ alone is an option only in the regime
$R'\!/R \ll \phi'$, where
$\Im\mathcal{C}/(\Delta^2\,|r|^2) \approx \phi'(\delta) + \phi'\,R'/R \to \phi'$.

\section{Implementation notes}

\begin{itemize}\setlength\itemsep{2pt}
  \item $r$ is stored on a detuning grid; off-grid $r(\delta\pm\Delta)$
        is evaluated by complex linear interpolation (real and imaginary
        parts independently).
  \item Scan $\delta \in \delta_\min + \Delta,\;\delta_\max - \Delta$ to
        keep both sidebands in-band; outside this range truncate.
  \item The overall prefactor $|E_0|^2 J_0(m) J_1(m)$ is absorbed into
        normalization; only the $\delta$- and $\Delta$-dependence of
        $\mathcal{C}$ matters.
\end{itemize}



\appendix




\newpage

\section{Fiber Bragg grating}
The same summary as in Ref. \cite{note_temp}.
\label{app:1dmodel}
\subsection{1D model: uniformly alternate between $n$ and $n+\delta n$}
The FBG can be described by a 1D model with alternating two layers ($n_1, a_1)$ and $n_2, a_2)$. 
Analytical expressions of reflection and transmission coefficients are well know and convenient for later investigation as described in Ref. \cite{Bendickson1996}.
Here is a brief summary of key equations:
The transfer matrix of two layers ($n_1, a_1)$ and $n_2, a_2)$ is expressed by
\begin{eqnarray}
    M_{\mathrm{cell}}=\left(\begin{array}{cc}
1 / t^* & r / t \\
r^* / t^* & 1 / t
\end{array}\right)
\label{eq:tmat}
\end{eqnarray}
where $\phi_1=n_1 k a_1$, $\phi_2=n_2 k a_2$, and
\begin{eqnarray}
\frac{r}{t}  =\frac{i}{2}\left(\frac{-n_1}{n_2}+\frac{n_2}{n_1}\right) \sin \left(\phi_1\right) e^{-i\phi_2}, \quad
\frac{1}{t}  =\left(\cos \left(\phi_1\right)+\frac{i}{2}\left(\frac{n_1}{n_2}+\frac{n_2}{n_1}\right) \sin \left(\phi_1\right)\right) e^{i\phi_2}
\end{eqnarray}
The dispersion relation of an infinite 1D chain is given by using a unit cell length $a=a_1+a_2$ as
\begin{eqnarray}
%\mathrm{Tr}[M_\mathrm{cell}]&=&2\cos\beta a\\
    %\rightarrow
    \quad \cos (\beta a)&=&\cos \left(k n_1 a_1\right) \cos \left(k n_2 a_2\right)-\frac{1}{2}\left(\frac{n_1}{n_2}+\frac{n_2}{n_1}\right) \sin \left(k n_1 a_1\right) \sin \left(k n_2 a_2\right)
\end{eqnarray}
Then, $N$-periodic system is written by 
\begin{eqnarray}
    M_N=M_\mathrm{cell}^N = M_\mathrm{cell}\frac{\sin(N\beta a)}{\sin(\beta a)}-I\frac{\sin((N-1)\beta a)}{\sin(\beta a)}.
\end{eqnarray}
Thus, transmission, reflection coefficienta and phase are written by \cite{Bendickson1996}
\begin{eqnarray}
    \frac{r_N}{t_N}=\frac{r}{t}\frac{\sin(N\beta a)}{\sin(\beta a)}, \quad
    \frac{1}{t_N}=\frac{1}{t}\frac{\sin(N\beta a)}{\sin(\beta a)}-\frac{\sin((N-1)\beta a)}{\sin(\beta a)}
\end{eqnarray}
As shown in Fig. \ref{fig:dn_dep}, the larger bandgap ($\propto \delta n/n$) leads to the smaller $\partial \phi_\mathrm{FBG}/\partial T$.
%Note that the quarter-wave stack model is convenient to derive a compact form of $R/T$ and phase as shown in Ref. \cite{Bendickson1996}.
\begin{figure}[!h]  
        \centering
        \includegraphics[width=1\linewidth]{fig/phase_21_20_17.png}
	\caption{Cell number dependence of uniform FBG with $\delta n/n=3.6\times 10^{-4}$: $N_\mathrm{cell}=5000, 10000, 15000$, corresponding to FBG length of $1.8$ mm, 3.6 mm and 5.3 mm, respectively. At the center of the bandgap, $d\phi_\mathrm{FBG}/d\nu = -0.025$ rad/GHz, independent of the cell number.}
    \label{fig:ncell_dep}
\end{figure}

\begin{figure}[!h]  
        \centering
        \includegraphics[width=1\linewidth]{fig/phase_22_17_23_25_26_24.png}
	\caption{Modulation $\delta n$ dependence with $N_\mathrm{cell}=15000$: As increasing the modulation amplitude $\delta n$, the band gap increases owing to $\Delta\nu_\mathrm{gap}/\nu=\delta n/n$. As a result, the slope the phase shift approaches to zero with a asymptotic form of $\frac{\partial \phi_\mathrm{FBG}}{\partial \nu} = \frac{A}{1+(\delta n/n)/B}$ where $A=-11.2(2.0), B=9.5(1.7)\times 10^{-7}$.}
    \label{fig:dn_dep}
\end{figure}
\vspace{-7mm}
\subsection{1D model: gaussian envelope of $\delta n$}
Using the transfer matrix of Eq. (\ref{eq:tmat}) with spatially dependent $\delta n$ (gaussian distribution with $\sigma = 5.9$ mm, N=17500), we numerically compute the $\tilde{M}_N = \prod_i^N M_\mathrm{cell}(\delta n_i)$ as shown in Fig. \ref{fig:fbgtuningGaussian}



\bibliography{mybib} 

\end{document}
